
    \expectedproblemsection{Mathematical Intuition Problems}

    \begin{expectedproblem}
        A geodesic triangle is drawn on a surface of constant Gaussian
        curvature. Predict whether its angle sum is less than, equal to, or
        greater than $\pi$ in each case.
        \begin{problemsubparts}
            \item negative curvature;
            \item zero curvature;
            \item positive curvature.
        \end{problemsubparts}
    \end{expectedproblem}

    \begin{solution}
        By the Gauss--Bonnet relation
        \[
            \alpha+\beta+\gamma-\pi
            =
            K\area(\triangle ABC),
        \]
        the sign of the angle excess is the sign of $K$. Therefore,
        \[
            \boxed{
                \begin{aligned}
                    \text{(a)}\;&\text{less than }\pi,\\
                    \text{(b)}\;&\text{equal to }\pi,\\
                    \text{(c)}\;&\text{greater than }\pi.
                \end{aligned}
            }
        \]
    \end{solution}

    \begin{expectedproblem}[1996 Putnam]
        Two circles have radii $1$ and $3$, and the distance between their
        centers is $10$. A point $X$ moves on the first circle and a point $Y$
        moves on the second. Determine the locus of the midpoint $M$ of
        $\overline{XY}$.
    \end{expectedproblem}

    \begin{solution}
        Let the centers of the two circles be $A$ and $B$, and let $O$ be the
        midpoint of $\overline{AB}$. Write
        \[
            \overrightarrow{AX}=\mathbf{u},
            \qquad
            \overrightarrow{BY}=\mathbf{v},
        \]
        where
        \[
            \|\mathbf{u}\|=1,
            \qquad
            \|\mathbf{v}\|=3.
        \]
        Since $M$ is the midpoint of $X$ and $Y$,
        \[
            \overrightarrow{OM}
            =
            \frac{\mathbf{u}+\mathbf{v}}{2}.
        \]
        The possible lengths of $\mathbf{u}+\mathbf{v}$ range from
        \[
            3-1=2
            \qquad\text{to}\qquad
            3+1=4.
        \]
        Every intermediate length and every direction occur by varying the
        angle between $\mathbf{u}$ and $\mathbf{v}$ and rotating both vectors.
        Therefore the locus is the closed annulus
        \[
            \boxed{1\leq OM\leq2}.
        \]
    \end{solution}

    \begin{expectedproblem}[1996 Putnam]
        Two squares have total area $1$. Their sides, and the sides of a
        containing rectangle, must remain parallel, and the interiors of the
        two squares may not overlap. Determine the least number $A$ s.t.
        every such pair of squares can be packed into a rectangle of area at
        most $A$.
    \end{expectedproblem}

    \begin{solution}
        Let the side lengths be $x\geq y>0$, so
        \[
            x^2+y^2=1.
        \]
        Since two axis-parallel squares with disjoint interiors must be
        separated horizontally or vertically, any containing rectangle has one
        side at least $x+y$ and the other at least $x$. This bound is attained
        by placing the squares side by side. Thus the least area required for
        this pair is
        \[
            x(x+y).
        \]

        Put $t=y/x$, where $0<t\leq1$. Then
        \[
            x^2=\frac{1}{1+t^2},
        \]
        and the required area is
        \[
            \frac{1+t}{1+t^2}.
        \]
        Differentiation shows that this expression is maximized at
        \[
            t=\sqrt2-1.
        \]
        Its maximum value is
        \[
            \boxed{\frac{1+\sqrt2}{2}}.
        \]
    \end{solution}

    \begin{expectedproblem}[1994 CSAT]
        A cube is intersected by a plane. Among the following shapes, determine
        how many can occur as the cross-section:
        \begin{problemchoices}
            \item a triangle;
            \item a nonsquare rectangle;
            \item a nonsquare rhombus;
            \item a pentagon;
            \item a hexagon.
        \end{problemchoices}
    \end{expectedproblem}

    \begin{solution}
        Work with the cube
        \[
            0\leq x,y,z\leq1.
        \]
        Each of the five shapes occurs.

        The plane
        \[
            x+y+z=\frac12
        \]
        produces a triangle. The plane
        \[
            z=\frac{x}{2}
        \]
        produces a rectangle whose side lengths are $1$ and $\sqrt5/2$, so it
        is not a square.

        The plane
        \[
            4z=x+y+1
        \]
        produces a parallelogram with adjacent side vectors
        \[
            \left(1,0,\frac14\right)
            \qquad\text{and}\qquad
            \left(0,1,\frac14\right).
        \]
        These vectors have equal lengths but nonzero dot product, so the
        cross-section is a nonsquare rhombus.

        The plane
        \[
            x+y+2z=\frac{11}{10}
        \]
        meets five edges of the cube, at
        \[
            \left(\frac1{10},1,0\right),
            \left(1,\frac1{10},0\right),
            \left(1,0,\frac1{20}\right),
            \left(0,0,\frac{11}{20}\right),
            \left(0,1,\frac1{20}\right),
        \]
        and therefore produces a pentagon. Finally,
        \[
            x+y+z=\frac32
        \]
        produces a hexagon. Thus all five listed shapes are possible, and the
        required number is
        \[
            \boxed{5}.
        \]
    \end{solution}
